\chapter{The fit check}
\label{ch:fit}

\section*{What this chapter is for}

Some reasons you will not get an offer have nothing to do with how well you
prepare, and they are knowable in advance. This chapter is a twenty-minute
check you run \emph{before} applying, and it exists because the alternative is
finding out in week four.

\section{Why this goes first}

Preparation is expensive and it is spent up front. A week of study, a portfolio
artefact, research on four people, a rehearsal. That investment is rational
when the process can be won and irrational when it cannot, and the things that
decide which are mostly visible on day one.

\judgement{The instinct, when a role is attractive, is to treat structural
mismatches as negotiable and get on with preparing. Sometimes they are
negotiable. The point of the check is not to refuse those roles; it is to price
them honestly, so that a week spent is a week you decided to spend.}

\section{The check}

Six questions. Answer each with what the evidence says, not with what you hope.

\begin{kittable}{@{}lXl@{}}
\textbf{Question} & \textbf{Where the answer is} & \textbf{If it is bad} \\
\midrule
Work model & the posting's metadata field, not its prose & usually fixed \\
Location & whether the role is advertised in several named cities & usually fixed \\
Band ceiling & published range, or the level in a public ladder & fixed at the top \\
Level & does the posting describe the level you are applying at & sometimes movable \\
Visa, notice, start date & your own facts & fixed \\
The segment & Chapter~\ref{ch:segment} & you can change which role you pursue \\
\end{kittable}

\textbf{Work model} is the one most often misread, and the misreading is
predictable: the body text says \emph{flexible} and the metadata field says
\emph{on-site}. The metadata is the one the system was configured with. Where
they disagree, believe the field and ask the question in the screen.

\textbf{The band ceiling} is a hard edge more often than people expect. Where a
company publishes a ladder, the advertised range usually maps to specific rows
in it, and the rows above are frequently not filled from outside. If your
expectation sits above the top of the range, the process can go perfectly and
still not produce an offer you would accept. That is worth knowing in week
zero.

\textbf{Several named cities} in one posting almost always means a team
assignment rather than a remote role, whatever the body text implies.

\section{The go / no-go}

\begin{kitmethod}
Run the six, then classify:

\textbf{Green} \dash{} nothing structural is in the way. Prepare fully.

\textbf{Amber} \dash{} one structural mismatch that is plausibly negotiable.
Prepare, and \emph{raise it in the screen}. Not at the offer stage: in the
screen, where it costs a conversation instead of a month. If it cannot be
raised early without ending the process, it is not amber, it is red.

\textbf{Red} \dash{} a mismatch that is fixed, and that you would not accept.
Apply if you want the practice, and know that is what you are buying. Do not
spend the week.
\end{kitmethod}

\begin{kitnote}
An amber that you decide to carry silently is a red with a longer timeline. The
whole value of naming it early is that you find out whether it moves while the
cost of the answer is still small.
\end{kitnote}

\section{What this check cannot see}

Two things, and they belong in Chapter~\ref{ch:classes} rather than here
because they are about you rather than about the role.

The first is your own history as a reader will see it: the shape of your last
few roles, the gaps, the lengths. That is not a fit question about this
company; it is a property of your evidence, and it travels with you to every
process.

The second is whether you have demonstrated scope. A role can be a perfect
structural fit and still be a level above where your evidence currently
reaches.

Neither is a reason not to apply. Both are reasons to know which objection you
are likely to meet, so that the answer is prepared rather than improvised.

\begin{drill}
Fill in \code{templates/fit-check.md} for the role you are most serious about,
including the amber items and how you would raise each in a screening call.
Time it: it should take twenty minutes. If it takes an hour, the extra forty
minutes were spent arguing yourself out of something the evidence already said.
\end{drill}

\section*{What this chapter does not claim}

It does not claim structural mismatches never move. They do \dash{} work models
get negotiated, levels get adjusted, bands get exceeded for specific people.
The claim is that they move rarely enough that it is worth pricing them as
fixed, and early enough that you find out while it is cheap.

The traffic-light classification is a decision aid, not a measurement. It has
no evidence behind it beyond being a structure for a decision you are making
with or without it.
